\chapter{Variations, extrémums et optimisation}
\section*{Exercice 16.1 -- Lire des variations}
Le minimum vaut $-1$, atteint en $1$ ; le maximum vaut $8$, atteint en $-2$. Les deux autres entrées appartiennent à des intervalles de variation différents ; ces seules informations ne permettent pas de comparer leurs images.

\section*{Exercice 16.2 -- Borne non atteinte}
Elle vérifie $0<f(x)<4$. Les bornes ne sont pas atteintes et la fonction est strictement croissante. Elle n'a donc ni minimum ni maximum sur cet intervalle ouvert.

\section*{Exercice 16.3 -- Optimiser}
Ses côtés sont $x$ et $12-x$, avec $0<x<12$. L'aire vaut $x(12-x)=36-(x-6)^2\leq36$. L'égalité est atteinte pour $x=6$, admissible. Le carré de côté $6$ réalise l'aire maximale $36\unite{m^2}$.

\section*{Exercice 16.4 -- Dichotomie}
Après $n$ étapes, l'amplitude vaut $2^{-n}$. Or $2^6=64<100$ et $2^7=128\geq100$. Il faut au moins $7$ étapes ; l'amplitude vaut alors $1/128\simeq0{,}0078125$.

\section*{Exercice 16.5 -- Grille et preuve}
Les valeurs sont $0{,}91$, $0{,}96$, $0{,}51$ : la grille retient $0{,}96$ en $0{,}5$. Mais $f(x)\leq1$, avec égalité en $0{,}3$, donc le maximum global vaut $1$. Une grille peut manquer le point optimal.

\section*{Exercice 16.6 -- Longueur approchée}
Avec un segment, la longueur vaut $\sqrt{1^2+1^2}=\sqrt2$. Avec deux segments, les différences de coordonnées sont $(1/2,1/4)$ puis $(1/2,3/4)$, d'où $L_2=(\sqrt5+\sqrt{13})/4\simeq1{,}4604$. Ce sont des longueurs polygonales, pas une formule exacte de la longueur de la parabole.
